\subsubsection{Modelo de negocio}
\label{sec:business-model}

El modelo de negocio del proyecto se sustenta en dos fuentes principales de ingresos. 
La primera proviene de la venta de gemas, que los usuarios adquieren para comprar cosméticos para sus personajes. 
La segunda fuente corresponde
a la comercialización de zonas para creadores de mundos, que mediante una suscripción mensual pueden alojar sus mundos en la plataforma.

\vspace{0.25cm}

Respecto a los costos, los principales corresponden al mantenimiento de los servidores, que incluyen el cómputo de los
mundos y el almacenamiento en los \textit{buckets} de los objetos, texturas y cosméticos asociados. Además, se requieren
inversiones en marketing para promocionar el juego y atraer nuevos usuarios. Finalmente, deben considerarse los costos de los 
desarrolladores, especialmente ante un escenario de sostenibilidad a largo plazo.

\vspace{0.25cm}

Por otro lado, se contempla la transferencia a los creadores de un pequeño porcentaje de las ganancias asociadas a la
venta de cosméticos. Sin embargo, esta funcionalidad no se ha implementado debido a su complejidad técnica y a que
queda fuera del alcance del MVP. Por ahora, únicamente existe un control de cuentas para registrar sus ganancias, por
lo que los creadores no pueden recibirlas de manera automática. 

\vspace{0.25cm}

Una posible estrategia de monetización futura podría ser la colaboración con marcas para crear mundos temáticos. Esto 
generaría ingresos por la asociación y, además, serviría como una estrategia de marketing para atraer nuevos usuarios 
a la plataforma.
